% Section 5 : Tests
% Chiffres et sources : resultats/chiffres_rapport.md (fin du fichier)

\section{Tests}
\label{sec:tests}

Notre objectif est simple : afficher exactement les mêmes lignes que
\texttt{grep -E}. Nous avons donc comparé notre programme à
\texttt{grep}, sur de vrais livres et sur des regex choisies pour piéger
chaque partie du code.

\paragraph{Les textes.}
Ce sont neuf livres du projet Gutenberg, en UTF-8 (tableau~\ref{tab:textes}).
Le script \texttt{tests/telecharger\_textes.sh} télécharge les sept
premiers ; \texttt{text1.txt} et \texttt{text2.txt} étaient déjà dans
\texttt{data/}. Le script retire les \texttt{\textbackslash r} : les
fichiers Gutenberg sont souvent en fin de ligne Windows, et
\texttt{grep} compterait le \texttt{\textbackslash r} comme un caractère
de la ligne. Il vérifie aussi que chaque fichier est bien en UTF-8, et
que l'ensemble reste sous 6~Mo (5\,333\,057 octets au total). Le livre
56667 commence par un BOM UTF-8 (U+FEFF). Nous l'avons gardé : c'est un
caractère ordinaire pour \texttt{grep} comme pour nous, et le banc le
confirme.

\begin{table}[htbp]
\centering
\small
\begin{tabular}{|l|r|r|l|}
\hline
Titre, auteur & N\textsuperscript{o} Gutenberg & Taille (octets) & Langue \\
\hline
\emph{The Yellow Wallpaper}, C.~P.~Gilman & 1952 & 51\,344 & anglais \\
\emph{Metamorphosis}, F.~Kafka & 5200 & 139\,673 & anglais \\
\emph{Alice's Adventures in Wonderland}, L.~Carroll & 11 & 170\,549 & anglais \\
\emph{Candide, ou l'optimisme}, Voltaire & 4650 & 230\,056 & français \\
\emph{The Double Thirteen}, A.~Wynne (\texttt{text2.txt}) & 79867 & 433\,257 & anglais \\
\emph{Le tour du monde en quatre-vingts jours}, J.~Verne & 800 & 452\,511 & français \\
\emph{A History of Babylon}, L.~W.~King (exemple du sujet) & 56667 & 749\,807 & anglais \\
\emph{Pride and Prejudice}, J.~Austen (\texttt{text1.txt}) & 1342 & 757\,471 & anglais \\
\emph{Don Quixote}, M.~de Cervantes (trad.\ anglaise) & 996 & 2\,348\,389 & anglais \\
\hline
\end{tabular}
\caption{Les neuf textes de \texttt{data/}.}
\label{tab:textes}
\end{table}

\paragraph{Les 41 regex.}
Le fichier \texttt{tests/regex.txt} les range par catégorie : 8 littéraux
(courts, moyens, longs, et \texttt{anas}, qui a un bord), 5 alternatives
simples ou imbriquées, 6 étoiles simples ou imbriquées, 7 regex avec le
point (dont \texttt{.*} et \texttt{h.*e.*r}), 4 façons d'écrire le mot vide
(la regex vide, \texttt{a|}, \texttt{()}, \texttt{(a|)b}), 7 motifs avec
accents ou typographie (\texttt{é}, \texttt{déjà}, \texttt{Ça}, \texttt{’},
\texttt{n’t}\ldots), 2 motifs sans résultat et les 2 motifs du sujet
(\texttt{S(a|g|r)*on} et \texttt{(Merodach|Sargon).*wall}). Au total,
15 de ces regex sont littérales, c'est-à-dire de simples suites de
lettres : les 8 littéraux de la première catégorie, 5 motifs accentués
(\texttt{é}, \texttt{déjà}, \texttt{Ça}, \texttt{’}, \texttt{n’t}) et les
2 motifs sans résultat (\texttt{zzzq}, \texttt{Sargonzzz}). Les caractères
\texttt{+ ? [ ] \{ \} \^{} \$ \textbackslash} n'y figurent jamais : ce
sont des opérateurs pour \texttt{grep -E}, mais de simples lettres pour
notre analyseur, qui ne connaît que l'ERE simplifiée du sujet. Un écart
sur ces caractères ne dirait rien de nos algorithmes.

\paragraph{Le banc de vérification.}
\texttt{tests/verifier.sh} lance chaque regex sur chaque texte dans cinq
configurations : les options par défaut, le matcher naïf, le matcher une
passe, KMP forcé (pour les 15 regex littérales seulement) et l'option
\texttt{-n} face à \texttt{grep -n -E}. Pour chaque cas, il compare
octet par octet notre sortie standard à celle de \texttt{grep -E}, et
compare aussi les codes de sortie. Cela fait $41 \times 9 = 369$ cas par
configuration, et 135 pour KMP : 1\,611 comparaisons en tout,
\textbf{0~écart}. Le banc tourne en 100~s environ et s'arrête avec le code~1
au moindre écart. Aucun de nos textes ne contient de caractère au-delà de
U+FFFF (un emoji, par exemple) : le banc ne couvre donc pas ce cas, où
notre point lit un demi-caractère. Cet écart est décrit avec les limites
du programme.

\paragraph{Tester le banc.}
Un banc qui ne trouve jamais rien ne prouve rien. Nous l'avons donc
lancé sur des versions volontairement fausses du code. Un KMP faux, qui
repart de 0 après un échec au lieu de reprendre en $q = \pi[q-1]$,
passait tous les tests tant qu'aucun motif du banc n'avait de bord.
Nous avons alors ajouté \texttt{anas} : son préfixe \texttt{ana} a le
bord \texttt{a}, et le KMP faux manque le mot \emph{ananas}, une ligne
dans \emph{Candide} et une dans le livre de Verne. Avec deux erreurs à la
fois (ce KMP faux et une table de transitions fausse), le banc trouve
515~écarts dans les cinq configurations.

\paragraph{Le francophone.}
\emph{Candide} et le livre de Verne sont là pour les accents. Lancé sur
une machine sans locale UTF-8, le banc donnait 146 faux écarts, sur des
motifs comme \texttt{é} ou \texttt{’}. La cause n'était pas dans notre code : sans
locale UTF-8, la JVM décode les arguments de la ligne de commande en
ASCII, et le motif arrive abîmé (caractère U+FFFD) avant même notre
analyseur. On ne peut plus changer ce décodage une fois la JVM lancée.
La correction : le banc exporte une locale UTF-8 (\texttt{C.UTF-8}, sinon
\texttt{en\_US.UTF-8}) pour \texttt{java} comme pour \texttt{grep}, et
le programme affiche un avertissement sur la sortie d'erreur s'il
détecte ce cas.

\paragraph{Test unitaire de KMP.}
\texttt{tests/TestKmp.java} compare notre KMP à \texttt{String.contains}
sur 20\,012 cas : 12 cas choisis à la main (motifs à bord comme
\texttt{aab} dans \texttt{aaab}, motif d'une lettre, motif plus long que
la ligne, ligne vide\ldots) et 20\,000 couples (motif, ligne) tirés au
hasard sur l'alphabet $\{a, b\}$, avec une graine fixe pour que chaque
lancement teste les mêmes cas. Tous réussissent.
